\subsection{Underwater image formation, turbidity and simulation}
\label{sec:lit-image-formation}

The light attenuation physics of underwater imaging is described by the Jaffe--McGlamery model, which is composed of three components at the sensor. These are the direct light attenuated along the line of sight, the forward-scattered light that blurs the direct image, and the backscattered veiling light that adds a range dependent haze \cite{mcglamery1980computer,jaffe1990computer}. Most applications uses a simplified form that keeps only the attenuation and veiling components, $I = J e^{-\beta d} + B(1 - e^{-\beta d})$. This simplification though has known limits. Akkaynak and Treibitz showed that the standard form mixes up two physically distinct attenuation coefficients, with consequences for colour reconstruction \cite{akkaynak2018revised}. Further, Ismiroglou et al.'s recent tank experiments put the forward-scatter (blur) term back into the simplified model. Their results show that in highly turbid water, images blur even at short distance to the camera \cite{ismiroglou2025seaing}. Part~B of this work adopts the simplified model, measures its $\beta$ from the target itself in the frame, then tests the blur question at the optical depths reached.

Marker detection under turbidity has also been approached with learning methods. Risholm et al. trained a marker detector with predictive uncertainty, which resulted in 100\% detection rate across attenuation lengths where the standard ArUco library's detection rate collapses, with the median translation error growing from 2.6~cm at low turbidity to 10.5~cm in higher turbidity \cite{risholm2021underwater}. DeepTurbid's approach replaces the four-corner detection with 68 learning keypoints per marker and trains on synthetically degraded imagery. It reports 95.8\% recall against OpenCV's 9.9\% on its turbid test set \cite{meng2026deepturbid}. Both results show that learning buys detection recall in murky water at a cost in compute and localisation precision, so the classical (PnP) pipeline remains the right baseline for this thesis. Notably, DeepTurbid's training data is generated by re-attenuating clear in-air imagery through the simplified image formation model. This procedure is used in Part~B, by degrading real underwater frames with geometric range truth and a measured attenuation coefficient.

Simulation platform occupy the other end of the fidelity spectrum. The standard robotics simulator like Gazebo model ridgid-body dynamics, buoyancy and hydrodynamic drag well enough to test the full control stack. However, its renderer model neither light attenuation nor veiling light, so simulated marker detection is optimistic. Physically grounded underwater image simulation exists such as OceanSim, but it does not accurately model vehicle and fluid dynamics, while Gazebo model the dynamics but not the optics. Currently, no open-sourced robotics platform provides both halves with validated fidelity, so this division motivates the two-part structure of this thesis.
